\begin{definition}[Natural-deduction derivation]
A \emph{natural-deduction derivation} is a finite rooted tree generated by the
rules displayed below.  Leaves may be marked assumptions.  An application of
an introduction or elimination rule may discharge the marked assumptions
indicated in brackets.  An assumption is \emph{open} iff no rule on the path
from that assumption to the root discharges it.

Discharge labels are metalinguistic markers drawn from \(\N\).  A consistent
injective renaming of the labels on assumptions and of the matching labels on
rules preserves a derivation and which assumptions it discharges.  Since a
derivation is finite, copies of derivations can always be given pairwise
disjoint label sets fresh from every label in a fixed receiving derivation.

An application of \(\to I\), \(\neg I\), or \(\mathsf{RAA}\) carries a
label \(i\) and discharges exactly the open occurrences marked \(i\) of
the bracketed assumption above its premise; an application of \(\lor E\)
carries two labels \(i,j\), as displayed, and discharges exactly the open
occurrences of the first disjunct marked \(i\) above its second premise
and the open occurrences of the second disjunct marked \(j\) above its
third premise.  Every label carried by such an application discharges at
least one occurrence.  Marking any set of open occurrences of one formula
with a common label fresh from every label in the derivation preserves
the derivation and the occurrences discharged by each of its rule
applications, since no rule application carries the fresh label; this
merging relabeling is distinct from the injective renamings above, and is
the operation used when a newly added root inference is to discharge
several open occurrences at once.  There is no repetition step: where an
informal presentation repeats a formula, the tree reuses the node already
present.
\end{definition}

The active calculus \(\ND_{\CPL}\) has the following rules.
\[
\begin{array}{c@{\qquad}c@{\qquad}c}
\Rule[\(\land I\)]{\varphi\quad\psi}{\varphi\land\psi}
&
\Rule[\(\land E_1\)]{\varphi\land\psi}{\varphi}
&
\Rule[\(\land E_2\)]{\varphi\land\psi}{\psi}
\\[2.2ex]
\Rule[\(\lor I_1\)]{\varphi}{\varphi\lor\psi}
&
\Rule[\(\lor I_2\)]{\psi}{\varphi\lor\psi}
&
\Rule[\(\to E\)]{\varphi\quad\varphi\to\psi}{\psi}
\end{array}
\]
\[
\Rule[\(\lor E\)]{
  \varphi\lor\psi\qquad
  \begin{array}{c}[\varphi]^i\\ \vdots\\ \chi\end{array}
  \qquad
  \begin{array}{c}[\psi]^j\\ \vdots\\ \chi\end{array}}
 {\chi}
\side{discharging all assumptions marked \(i\) and \(j\)}.
\]
\[
\begin{array}{c@{\qquad}c@{\qquad}c}
\Rule[\(\to I\)]{
  \begin{array}{c}[\varphi]^i\\ \vdots\\ \psi\end{array}}
 {\varphi\to\psi}
&
\Rule[\(\neg I\)]{
  \begin{array}{c}[\varphi]^i\\ \vdots\\ \bottom\end{array}}
 {\neg\varphi}
&
\Rule[\(\neg E\)]{\varphi\quad\neg\varphi}{\bottom}
\\[4ex]
\Rule[\(\bottom E\)]{\bottom}{\varphi}
&
\Rule[\(\mathsf{RAA}\)]{
  \begin{array}{c}[\neg\varphi]^i\\ \vdots\\ \bottom\end{array}}
 {\varphi}
&
\end{array}
\]
The rules \(\to I\), \(\neg I\), and \(\mathsf{RAA}\) discharge the
indicated assumptions.  The last rule is the explicitly classical rule.

\begin{definition}[Propositional derivability]
For \(\Gamma\subseteq\Fm_{\CPL}\), write
\(\Gamma\derives_{\ND_{\CPL}}\varphi\) iff there is a
\(\ND_{\CPL}\)-derivation of \(\varphi\) whose open assumptions all belong to
\(\Gamma\).  A \emph{proof} is such a derivation together with its stated set
of permitted open assumptions.
\end{definition}

\begin{example}[Commutativity of conjunction]
The following derivation proves
\(\derives_{\ND_{\CPL}}(P\land Q)\to(Q\land P)\):
\[
\begin{array}{rll}
1.&[P\land Q]^1&\text{assumption}\cr
2.&Q&\land E_2,1\cr
3.&P&\land E_1,1\cr
4.&Q\land P&\land I,2,3\cr
5.&(P\land Q)\to(Q\land P)&\to I,1\text{--}4.
\end{array}
\]
The superscript marks the assumption discharged at line 5.
\end{example}

\begin{example}[Excluded middle]
The formula \(P\lor\neg P\) has a classical proof:
\[
\begin{array}{rll}
1.&[\neg(P\lor\neg P)]^1&\text{assumption}\cr
2.&\quad[P]^2&\text{assumption}\cr
3.&\quad P\lor\neg P&\lor I_1,2\cr
4.&\quad\bottom&\neg E,1,3\cr
5.&\neg P&\neg I,2\text{--}4\cr
6.&P\lor\neg P&\lor I_2,5\cr
7.&\bottom&\neg E,1,6\cr
8.&P\lor\neg P&\mathsf{RAA},1\text{--}7.
\end{array}
\]
The proof uses \(\mathsf{RAA}\).  The calculus obtained by omitting
\(\mathsf{RAA}\) is intuitionistic natural deduction, made exact in
\cref{ch:intuitionistic-logic}; no derivation in that calculus concludes
\(P\lor\neg P\) from no open assumptions.
\end{example}

\begin{lemma}[Substitution for open assumptions]
\label{lem:nd-substitution}
If \(D\) derives \(\psi\) with an open assumption \(\varphi\), and \(E\)
derives \(\varphi\), then replacing each designated open occurrence of
\(\varphi\) in \(D\) by a copy of \(E\), after consistently renaming every
label in every copy to be globally fresh from \(D\) and from all other copies,
produces a derivation of \(\psi\).  Its open assumptions are those of \(D\)
other than the replaced occurrences, together with the correspondingly
renamed open assumptions of the inserted copies of \(E\).
\end{lemma}
\begin{proof}
First reserve for each designated occurrence a set of labels disjoint from
the labels of \(D\) and from the sets reserved for the other occurrences, and
consistently rename its copy of \(E\) into that set.  Strengthen the induction
invariant by requiring every inserted copy to use its reserved globally fresh
labels.  Induct on the height of \(D\).  If \(D\) is the designated
assumption, use its renamed copy of \(E\).  If it is another assumption, leave
it unchanged.  For an inference, apply the induction hypothesis to each
immediate subderivation and then reapply the same rule with the original
discharge labels of \(D\).  Consistent renaming preserves every discharge
internal to a copy of \(E\), while global freshness makes every label of that
copy different from every label on a rule of \(D\).  Hence no rule above a
graft point can accidentally discharge an open assumption of the inserted
copy, and the stated open-assumption accounting follows.
\end{proof}

